\bookmark[dest=\HyperLocalCurrentHref, level=1]{Conclusion}
\section{Conclusion} \label{sec6}
In this paper, we have presented a comprehensive survey of non-line-of-sight (NLOS) imaging, covering the foundational physical models and algorithms established through 2021 and the significant advances made from 2022 to 2026. The field has matured substantially during this period, progressing from laboratory demonstrations with minute-long acquisition times to real-time video of dynamic room-sized scenes.

\textbf{Progress on fundamental challenges.}
The three core challenges identified in the 2021 survey---ill-posedness with low SNR, limited resolution, and long data acquisition time---have all seen meaningful progress. On SNR, photon-efficient imaging~\cite{liuPhotonEfficientNLOS2022} and first-photon event stamping~\cite{liFastFirstPhoton2022} have demonstrated high-quality reconstruction at sub-photon-per-pixel levels. Resolution has been pushed to $\sim0.6\,\text{mm}$ via active wavefront shaping (UNCOVER~\cite{caohighresolutionnlos2022}), and instrument response function deconvolution further improves spatial resolution without hardware changes. Most dramatically, real-time 4\,fps NLOS video of room-sized dynamic scenes has been demonstrated~\cite{yeRealTimeNLOS2024}, and scan-free systems operating at up to 19\,fps have been reported~\cite{zhangRealTimeScanFree2024}, fundamentally changing the prospects for practical NLOS deployment.

\textbf{Deep learning advances.}
The deep learning component of NLOS imaging has seen a transformation in architecture diversity and sophistication. Beyond the U-Net and ResNet baselines of the 2021 era, the field has adopted transformers (NLOST~\cite{liNLOST2023}, TransiT~\cite{liTransiT2025}), Mamba/state-space models for video~\cite{liMambaTemporalConsistency2024}, graph neural networks~\cite{suDGNLOS2025}, and diffusion model priors~\cite{suModelGuidedDiffusion2024}. Neural implicit surface representations (NLOS-NeuS~\cite{fujimuraNLOSNeuS2023}) have extended beyond voxel volumes to enable smooth 3D surface reconstruction. Learnable physical priors~\cite{sunGeneralizable2025} and domain reduction strategies~\cite{shimDomainReduction2024} have improved generalization, while unsupervised methods~\cite{cuiVirtualScanning2024} and untrained network priors~\cite{wuUntrainedDecoder2022} have reduced dependence on large training datasets. Dedicated NLOS rendering tools~\cite{royoNLOSTransientRendering2022,plackFastDifferentiable2023} and real-world datasets such as NLOS-Track~\cite{wangPropagateCalibrate2023} have better supported the training and evaluation of these methods.

\textbf{Passive NLOS advances.}
Passive NLOS imaging has progressed from simple 2D reconstruction with occluders to real-time passive tracking of moving persons (PAC-Net~\cite{wangPropagateCalibrate2023}) and full-color 3D passive NLOS imaging with ordinary cameras~\cite{czajkowskiTwoEdge2024}. Hyperspectral and multispectral fusion (Hyper-NLOS~\cite{chenHyperNLOS2024}) and the exploitation of rough surfaces as relay surfaces~\cite{liRoughSurfaces2025} have broadened the environmental conditions and surface types compatible with passive NLOS sensing.

\textbf{New modalities and scenes.}
Beyond optical NLOS, mmWave radar NLOS reconstruction (HoloRadar~\cite{laiHoloRadar2025}) and ultrasound synthetic aperture NLOS imaging~\cite{ultrasoundNLOS2025} have demonstrated full 3D scene reconstruction in non-optical modalities, enabling operation in darkness, through walls, and under conditions where optical methods fail. New NLOS application scenes have also emerged: human pose estimation~\cite{xiaoHumanPose2026}, polarization-based scanning-free single-pixel imaging~\cite{zhouPolarizationSinglePixel2026}, and higher-order bounce imaging (virtual mirrors~\cite{royoVirtualMirrors2023}) have significantly expanded the scope of NLOS sensing.

\textbf{Remaining challenges and future directions.}
Despite remarkable progress, several important challenges remain open. (1) \textit{Outdoor and strong ambient light operation:} most existing systems require controlled or low-ambient-light conditions. Coherent sensing approaches~\cite{huangCoherentNLOS2024} and temporal gating designs~\cite{riccardoFastGatedSPAD2022} are promising, but robust outdoor operation at long range remains difficult. (2) \textit{Large-scale real training data:} the gap between synthetic and real-world NLOS measurements continues to limit generalization. Automated data collection systems, more realistic rendering, and domain adaptation methods are needed. (3) \textit{Non-Lambertian and complex relay surfaces:} most algorithms still assume diffuse relay walls; exploiting specular~\cite{liRoughSurfaces2025} and non-planar surfaces~\cite{guNonPlanar2023} in a unified framework remains an open problem. (4) \textit{End-to-end system integration:} bridging the gap between research demonstrations and practical deployed systems---handling calibration, robustness, and computational cost---is a critical next step. (5) \textit{Semantic NLOS understanding:} going beyond shape reconstruction to semantic tasks such as human pose estimation~\cite{xiaoHumanPose2026}, action recognition, and object identification in hidden scenes represents an important emerging frontier.

In summary, NLOS imaging has transitioned from a laboratory curiosity to a technology on the verge of practical deployment across multiple modalities and application domains. Continued progress on the challenges outlined above will be essential for realizing the transformative potential of seeing around corners.
